% !TEX program = pdflatex
\documentclass[3p,11pt]{elsarticle}
\usepackage{amsmath,amssymb}
\usepackage{array}
\usepackage{booktabs}
\usepackage[tableposition=top]{caption}
\usepackage{enumitem}
\usepackage{tikz}
\usepackage{pgfplots}
\pgfplotsset{compat=1.18}
\usepackage{url}
\urlstyle{same}
\usepackage[hidelinks]{hyperref}
\emergencystretch=3em
% first-page footer: the page number alone
\makeatletter
\def\ps@pprintTitle{\let\@oddhead\@empty\let\@evenhead\@empty
  \def\@oddfoot{\hfil\thepage\hfil}\let\@evenfoot\@oddfoot}
\makeatother

% ---- notation ---------------------------------------------------------------
\newcommand{\um}{\,$\mu$m}
\newcommand{\kms}{km\,s$^{-1}$}
\newcommand{\flux}{erg\,cm$^{-2}$\,s$^{-1}$}
\newcommand{\mA}{m\AA}
\newcommand{\sci}[2]{#1\times10^{#2}}
\newcommand{\ion}[2]{#1\,\textsc{#2}}
\newcommand{\nuc}[2]{$^{#1}$#2}
\newcommand{\lxl}{$L_X/L_\mathrm{bol}$}

\begin{document}

\begin{frontmatter}

\title{A search for waste heat from Dyson spheres around nearby stars in archival infrared surveys}

\author{Charles C. Norton}

\begin{abstract}
A structure that intercepts a star's light re-radiates it as heat, so a Dyson sphere
shows as infrared emission beyond what the star alone gives. We test nearby stars for
that emission in ten ways with archival surveys. None of the $279$ optically dark Gaia
DR3 astrometric companions of $0.5\,M_\odot$ or more within $200$\,pc re-radiates more
than $6$ per cent of a main-sequence star's light at $300$\,K (median bound $0.3$ per
cent), nor do the $191$ companions of $0.3\,M_\odot$ or more at $3$ to $30$\,au implied by
Hipparcos-Gaia proper-motion anomalies within $50$\,pc (median bound $0.39$ per cent).
Stacked against matched field stars, the WISE colours of habitable-zone planet hosts
within $100$\,pc limit the mean fraction of their light intercepted by $300$\,K swarms to
$\sci{1.5}{-5}$. No bright, cold, moving AllWISE source without a Gaia star is an
enclosed star, no star within $50$\,pc has a $30$ to $100$\,K shell in the AKARI
far-infrared survey, none of $203$ stars within $15$\,pc brightened at $4.6$\um\ over the
decade of NEOWISE by more than about $9$ per cent, and no star that vanished from the
Palomar plates reappears with the $12$\um\ emission an enclosure requires. SPHEREx
spectra trace the excess of earlier Dyson-sphere candidates to background galaxies and a
radio source, and put the excesses of $37$ extreme mid-infrared stars at $486$ to
$2133$\,K, the temperatures of warm dust and cool companions, none below $400$\,K within
the $5$\um\ reach of the spectra. Each test's reach is stated: the enclosed-star search
finds a $300$\,K shell around a star of $0.01\,L_\odot$ to $210$\,pc, the AKARI match
finds a $30$ to $100$\,K shell holding $0.2$ per cent of a solar luminosity at $50$\,pc,
and a vanished source of $R=19$ enclosed at $200$ to $600$\,K would have been detected
by WISE.
\end{abstract}

\begin{keyword}
Technosignatures \sep Dyson spheres \sep Infrared: stars \sep Circumstellar matter \sep
Solar neighbourhood
\end{keyword}

\end{frontmatter}

% =============================================================================
\section{Introduction}\label{sec:intro}
% =============================================================================

Dyson proposed that a civilization using much of its star's output would surround the
star with collectors and re-radiate the energy as waste heat, and that the heat, at a
few hundred kelvin, would make the system a source of infrared emission that its
starlight does not explain~\cite{Dyson1960}. Searches have since used each infrared
survey in turn: IRAS for whole-sky limits on spheres at $100$ to $600$\,K~\cite{Carrigan2009},
WISE for galaxies whose stellar light is converted to mid-infrared
emission~\cite{Wright2014,Griffith2015}, and Gaia with 2MASS and WISE for partial spheres
around individual Milky Way stars~\cite{Suazo2022,Suazo2024,Contardo2024}. A detected
excess has natural rivals: debris discs, red companions, background galaxies and
young stars embedded in their clouds.

The nearest stars allow tests that need a known distance, mass or history. A companion
whose mass Gaia's astrometry measures but whose light it does not see is either a dim
remnant or a star enclosed; a structure built within a decade brightens its star in the
infrared from one epoch to the next; a star enclosed during the twentieth century would
vanish from the Palomar plates and reappear in the mid-infrared. This paper applies
these tests and the stacking and spectroscopic ones to archival surveys, and states each
result as a fraction of a star's light re-radiated at the temperature tested.

\subsection{Contributions}
\begin{enumerate}[leftmargin=*,itemsep=1pt]
\item Limits on the fraction of light re-radiated at $300$\,K by $279$ Gaia DR3 dark
companions and by $191$ companions inferred from proper-motion anomalies.
\item A stacked limit on warm swarms around habitable-zone planet hosts, and searches
for fully enclosed moving stars, cold far-infrared shells, decade-long brightening and
stars that vanished from the optical sky into the infrared.
\item SPHEREx $0.75$ to $5$\um\ spectra of published Dyson-sphere candidates against
matched controls, which place every excess in a natural class.
\end{enumerate}

% =============================================================================
\section{Tests}\label{sec:tests}
% =============================================================================

Table~\ref{tab:tests} lists the ten tests, their samples and results. Each compares a
measured infrared flux with what the star or its catalogued companions would give, and
expresses the difference as the fraction of the star's light that a black body at the
tested temperature would need to re-radiate.

\emph{Dark astrometric companions.} Gaia DR3's orbital solutions~\cite{Arenou2023} give
the mass of each companion and the flux ratio its orbit allows in the $G$ band. Within
$200$\,pc, $279$ companions of $0.5\,M_\odot$ or more are dark, contributing less than a
third of the light of a main-sequence star of their mass. Inside a complete sphere such
a star would re-radiate its main-sequence luminosity as a $300$\,K black body, whose
$12$ and $22$\um\ flux at the Gaia distance we compare with the AllWISE~\cite{Wright2010}
counterpart within $3''$.

\emph{Companions from proper-motion anomalies.} The difference between a star's
Hipparcos-Gaia long-term proper motion and its Gaia proper motion gives the least mass
of a companion anywhere between $3$ and $30$\,au~\cite{Kervella2022}. For the $191$
Hipparcos stars within $50$\,pc whose anomaly, at a signal-to-noise ratio of $5$,
requires a companion of $0.3\,M_\odot$ or more, a star of that mass inside a shell at
$200$, $300$ or $600$\,K would add predictable $12$ and $22$\um\ flux; the star's
$K_s-W3$ and $K_s-W4$ colours, against the $15$ sample stars nearest in $G-K_s$, bound
the fraction re-radiated. An excess is read again from the WISE All-Sky release, which
processed the same cryogenic $12$ and $22$\um\ exposures apart from AllWISE, against the
same neighbours.

\emph{Stacking at habitable-zone hosts.} A swarm intercepting a fraction $f$ of the
starlight at $300$\,K adds $12$ and $22$\um\ flux in proportion to $f$, so a population
of modest swarms shows as a mean colour excess. Hosts of planets receiving $0.32$ to
$1.78$ times the Earth's insolation~\cite{Kopparapu2013} within $100$\,pc in the NASA
Exoplanet Archive~\cite{Akeson2013}, and the other hosts as comparison, each receive up
to $25$ field stars matched in Gaia colour, magnitude and parallax. The inverse-variance
mean of the hosts' $W1-W3$ and $W1-W4$ excesses over their controls, against the same
statistic with a control standing for its host, limits the mean $f$.

\emph{Fully enclosed stars.} A shell that absorbs all of a star's light re-radiates it
at the shell's temperature: at $300$\,K around an M dwarf of $0.01\,L_\odot$ at $100$\,pc
it gives $W3\approx6.4$ with $W2-W3\approx5.7$, and Gaia sees nothing. Of the $36603$
AllWISE point sources with $W3<8$ at a signal-to-noise ratio above $10$ and
$W2-W3>1.5$, we keep those that CatWISE2020~\cite{Marocco2021} finds moving faster than
$40$\,mas\,yr$^{-1}$ at $5\sigma$, as a star within $100$\,pc would, with no Gaia DR3
source within $3''$, clean flags and unsaturated $W1$ and $W2$, and check each against
SIMBAD and for young stars, cores and clouds within $30'$. The $W3<8$ cut reaches a
$300$\,K shell around a star of $0.001$, $0.01$ and $0.1\,L_\odot$ to $67$, $210$ and
$670$\,pc, and shells at $200$ and $500$\,K to within $20$ per cent of those distances;
the motion cut keeps stars whose tangential velocity exceeds $19$\,km\,s$^{-1}$ at
$100$\,pc and $40$\,km\,s$^{-1}$ at $210$\,pc, about three quarters of the disc
population at the first and a third at the second.

\emph{Cold shells.} A structure at $30$ to $100$\,K peaks beyond the WISE bands, and a
main-sequence photosphere at $65$ to $160$\um\ lies far below the AKARI Far-Infrared
Surveyor's bright source catalogue~\cite{Murakami2007}. We match the $27394$ Gaia DR3
main-sequence stars within $50$\,pc to that catalogue within $20''$, and take the chance
count from the same match with every position moved by $3'$. The catalogue's $0.55$\,Jy
limit at $90$\um\ corresponds, for a shell at $30$ to $100$\,K, to $0.2$ to $0.6$ per
cent of a solar luminosity at $50$\,pc and $0.03$ to $0.09$ per cent at $20$\,pc; for
the M dwarfs that make up most of the sample, of $0.01\,L_\odot$ at $10$\,pc, it needs
$1$ to $2$ per cent of the star's light, so the test bounds complete or nearly complete
cold shells around the M dwarfs and partial ones around the solar-type stars.

\emph{Decade trends.} A structure assembled during 2014 to 2024 at $600$ to $1000$\,K
brightens its star year by year, more at $4.6$ than at $3.4$\um. For the $298$ Gaia DR3
stars within $15$\,pc with $W1$ of $8$ to $12.5$, the NEOWISE-R single
exposures~\cite{Mainzer2014} from December 2013 to August 2024 give $W1$ and $W2$ light
curves differential against field stars on the same frames, and a linear trend per
decade is judged against stars of similar brightness, since the photometry drifts with
brightness.

\emph{Vanished in the optical, present in the infrared.} A star enclosed after the 1950s
leaves the first Palomar survey's plates and returns at its enclosure's temperature: at
$300$\,K a star of $R=19$ appears near $W2=16.8$ and $W3=11.1$, at $200$\,K near
$W3=11.6$ and $W4=9.0$, and at $600$\,K near $W2=14.1$, against AllWISE's $5\sigma$
limits of about $17.5$, $11.6$ and $8.0$ in $W2$, $W3$ and $W4$, so an enclosure at any
of those temperatures of a source brighter than $R\approx19$, which is most of the
catalogue, would be detected in at least one band. The archive of
vanishing objects~\cite{Solano2022} built for the VASCO project~\cite{Villarroel2020}
lists Palomar sources absent from Pan-STARRS and Gaia; we keep those whose only
counterpart is infrared and red in $W1-W2$, away from the ecliptic, with a fixed
CatWISE2020 counterpart and no Legacy Surveys source~\cite{Dey2019} to $r=24$.

\emph{Breakthrough Listen targets.} The FGK stars of Breakthrough Listen's nearby-star
sample~\cite{Isaacson2017} observed in a set of 2016 L-band sessions, $170$ in all, are
tested for $K_s-W3$ or $K_s-W4$ above $0.25$\,mag at $5\sigma$ with clean photometry.

\emph{SPHEREx spectra.} SPHEREx images the sky in $102$ bands from $0.75$ to
$5$\um~\cite{Dore2014}. We measure each target, carried by proper motion to the epoch
of the images, by a point-spread fit after the zodiacal model in every image covering
it, and the same for two Gaia stars of matching colour, magnitude and distance that WISE
does not detect at $12$\um. Normalised at $1.0$ to $1.6$\um, a target's excess over its
controls at $2$ to $3$, $3$ to $4$ and $4$ to $5$\um\ is judged against one control
measured against the other. The targets are candidates of Project
Hephaistos~\cite{Suazo2024} and the $53$ FGK stars with the most extreme Gaia-predicted
$W2$ excesses of Contardo and Hogg~\cite{Contardo2024}, of which $51$ have controls. A black body
fitted from $1.6$ to $5$\um\ gives the temperature of each excess and the fraction of
the star's light it holds. The Hephaistos targets are ten candidates of that search and
one M dwarf with a warm excess, eleven in all.

\begin{table}[t]
\centering\small
\caption{The ten tests. $f$ is the fraction of a star's light intercepted and re-radiated
at the temperature tested.}
\label{tab:tests}
\begin{tabular}{>{\raggedright\arraybackslash}p{0.2\textwidth} >{\raggedright\arraybackslash}p{0.3\textwidth} >{\raggedright\arraybackslash}p{0.42\textwidth}}
\toprule
Test & Sample & Result \\
\midrule
Dark companions & $279$ Gaia DR3 orbital companions $\geq0.5\,M_\odot$, $<200$\,pc & none above $6$ per cent of its light at $300$\,K; median bound $0.3$ per cent \\
Proper-motion companions & $191$ stars $<50$\,pc, companion $\geq0.3\,M_\odot$ at $3$--$30$\,au & median bound $0.39$ per cent at $300$\,K; four $3\sigma$ excesses, all explained \\
Habitable-zone stacking & $69$ hosts ($W3$), $63$ ($W4$), $<100$\,pc, $25$ controls each & mean $f<\sci{1.5}{-5}$ ($12$\um) and $\sci{8.6}{-6}$ ($22$\um) at $300$\,K \\
Fully enclosed stars & $36603$ cold AllWISE sources & $72$ clean movers without Gaia stars, all among young stars, cores and clouds \\
Cold shells & $27394$ stars $<50$\,pc, AKARI FIS & $35$ matches against $26$ by chance; clouds and known young discs \\
Decade trends & $203$ stars $<15$\,pc, NEOWISE 2013--2024 & none brightens at $4.6$\um\ at $5\sigma$; bound about $9$ per cent \\
Vanished stars & $5893$ red infrared-only counterparts of vanishing Palomar sources & four red CatWISE sources left, none detected at $12$\um \\
Breakthrough Listen targets & $170$ FGK stars & three excesses: a debris disc, a blend, a catalogued $22$\um\ excess \\
SPHEREx, Hephaistos & $11$ candidates, $2$ controls each & background galaxies and a radio source \\
SPHEREx, Contardo and Hogg & $51$ stars, $2$ controls each & $37$ excesses at $486$--$2133$\,K, none below $400$\,K \\
\bottomrule
\end{tabular}
\end{table}

% =============================================================================
\section{Results}\label{sec:results}
% =============================================================================

\emph{Companions.} No dark astrometric companion shows the $12$ or $22$\um\ flux that a
star of its mass inside a $300$\,K sphere would give: the measured flux bounds the
re-radiated fraction below $6$ per cent for every one of the $279$ and below $0.3$ per
cent for half of them. The proper-motion companions give a median bound of $0.39$ per
cent at $300$\,K, and $154$ of $189$ with both bands are bounded below $1$ per cent.
Three of their four $3\sigma$ colour excesses are a diffraction spike from a neighbour,
the photosphere of an M0 star and the white-dwarf companion of 56~Per. The fourth,
$0.14$\,mag at $12$\um\ ($3.3\sigma$) in the F6~V star HD~59099, whose astrometry
(RUWE $11.9$) marks the companion, belongs to the AllWISE measurement: the All-Sky
release puts the star $0.11$\,mag fainter at $12$\um\ and $0.07$\,mag fainter at
$22$\um, where the two releases agree for the sample's stars to a median of
$0.002$\,mag with a scatter of $0.007$\,mag, and against the same neighbours its
All-Sky excess is $+0.02$\,mag at $12$\um\ and nothing at $22$\um.

\emph{Habitable-zone hosts.} The $69$ habitable-zone hosts with clean $12$\um\
photometry and controls show a mean $W1-W3$ excess of $+0.0004\pm0.0045$\,mag over their
controls, the $124$ other hosts $-0.0011\pm0.0036$\,mag, and controls standing for their
hosts $+0.0018\pm0.0006$\,mag; at $22$\um\ the $63$ habitable-zone hosts give
$-0.0005\pm0.0094$\,mag, the $88$ other hosts $-0.0060\pm0.0106$\,mag, and controls
standing for their hosts $+0.0123\pm0.0017$\,mag. The control null is offset from zero
in both bands, by $0.0018$\,mag at $12$\um\ and $0.012$\,mag at $22$\um: a control
measured against the mean of the others is systematically redder than that mean, which
the hosts share, so the hosts' excesses are read against the null and fall below it by
$0.3\sigma$ and $1.3\sigma$. A $300$\,K swarm therefore intercepts on average less than
$\sci{1.5}{-5}$ of a habitable-zone host's light ($\sci{8.6}{-6}$ at $22$\um). The one
host above $3\sigma$, TOI-2257, has a galaxy $7.3''$ away.

\emph{Enclosed stars and cold shells.} Of the $36603$ cold AllWISE sources, $7817$ move
at $5\sigma$, $1491$ of these have no Gaia source, and $72$ have clean flags below
saturation; every one of the $72$ lies among young stars, dense cores or clouds, where
variability and nebulosity bias CatWISE motions. Within $50$\,pc, $35$ stars have an
AKARI $90$\um\ source within $20''$ against $26$ when the positions are moved; the
matches within $8''$ are alignments with clouds in fields of $43$ to $349$ sources per
square degree (V854~Ori, V1497~Cyg, HD~324218) or known young discs (HD~98800 in the TW
Hya association, and HD~164249~A in the $\beta$~Pictoris group at $\sci{7}{-4}$ of its
light).

\emph{Decade trends and vanished stars.} Of the $203$ stars within $15$\,pc with both
NEOWISE light curves, none brightens at $4.6$\um\ at $5\sigma$, the most extreme trend
being $-2.8\sigma$, and the median slope error of $0.018$\,mag per decade bounds any gain
at about $9$ per cent of the $4.6$\um\ flux. Of the $5893$ red infrared-only
counterparts of vanished Palomar sources, those near the ecliptic are asteroids caught
in single NEOWISE exposures and those detected in one exposure without a CatWISE source
are transients; four red CatWISE counterparts remain, and none is detected at $12$\um,
where an enclosure's energy would appear, consistent with brown dwarfs and obscured
quasars aligned by chance.

\emph{Breakthrough Listen targets.} Three of the $170$ stars exceed the colour criterion:
HR~8799, whose debris disc is known, a blended binary, and a star with a catalogued
$22$\um\ excess.

\emph{SPHEREx spectra.} The SPHEREx spectra of the Project Hephaistos candidates show
excesses only where a background galaxy (candidates D and E) or a radio source
(candidate G) lies within the WISE beam. JWST imaging of candidates D and E has since
resolved the galaxies, within about $1''$ of each star, as the source of their
mid-infrared excess~\cite{Zackrisson2026}, which SPHEREx's $6''$ pixels could attribute
only by the spectral shape. For the Contardo and Hogg stars, the ratio to
the controls shows a $4$ to $5$\um\ excess above $3\sigma$, with a systematic floor of
$0.063$ from $50$ control pairs, in $37$ of $51$ and a $3$ to $4$\um\ excess in $36$,
against $1$ and $0$ control pairs. A black body puts every excess between $486$ and
$2133$\,K holding $0.6$ to $7$ per cent of the star's light: $22$ at $400$ to $1000$\,K,
as warm dust, and nine above $1700$\,K, as cool companions or blends; none lies below
$400$\,K, and no $3.3$\um\ aromatic band reaches $2\sigma$. The absence of cooler
components is a property of the spectra's reach as much as of the stars: a $300$\,K
body emits at $5$\um\ a hundredth of what it emits at its $10$\um\ peak, so the fit from
$1.6$ to $5$\um\ is sensitive to components above about $400$\,K, and a cooler
component that the $12$ and $22$\um\ WISE excesses of these stars might hold is not
excluded. In the $41$ stars that the
Legacy Surveys cover, the excess sits in the star's own deblended source, whose $W2/W1$
ratio exceeds its controls' by up to $0.85$, and no source within $8''$ carries more than
$6$ per cent of it; the unWISE $W2$ centroid of every star lies within $3\sigma$ of its
$W1$ centroid against the controls' $0.12''$ scatter, which leaves only a red source
within about $0.5''$. The eleven with the largest excesses keep them at every NEOWISE
visit from 2014 to 2024, with a visit-to-visit scatter of $0.02$ to $0.06$\,mag in
$W1-W2$.

% =============================================================================
\section{Related work}\label{sec:related}
% =============================================================================

\emph{Waste-heat searches.} Carrigan's IRAS search set whole-sky limits on spheres at
$100$ to $600$\,K~\cite{Carrigan2009}, and the $\hat{G}$ search found no galaxy
re-radiating most of its starlight~\cite{Wright2014,Griffith2015}. Project Hephaistos
selected partial-sphere candidates among Milky Way stars with Gaia, 2MASS and
WISE~\cite{Suazo2022,Suazo2024}, and radio and optical images later showed background
sources near several of them~\cite{Ren2024}, as the SPHEREx spectra here do. Three
further studies of those candidates have appeared: archival diagnostics find
background contamination in candidates B and C and suggestive signs in five
more~\cite{Ren2026}, their stellar parameters are those of ordinary M dwarfs whose
infrared colours resemble transitional discs without the signs of
youth~\cite{Korn2026}, and JWST resolves background galaxies within about $1''$ of
candidates D and E~\cite{Zackrisson2026}, the two whose SPHEREx spectra attribute the
excess to a galaxy here. Contardo and
Hogg modelled the mid-infrared colours of five million FGK stars from their Gaia spectra
and published the most extreme excesses~\cite{Contardo2024}.

\emph{Companions and vanishing sources.} Kervella et al.\ derived companion masses from
the Hipparcos-Gaia proper-motion anomaly~\cite{Kervella2022}, and the VASCO project
catalogued Palomar sources missing from modern surveys~\cite{Villarroel2020,Solano2022};
we find no earlier use of either as a test for enclosed stars.

% =============================================================================
\section{Conclusion}\label{sec:conclusion}
% =============================================================================

No nearby star shows the infrared emission of a Dyson sphere in these tests. Optically
dark companions of known mass re-radiate at most a few tenths of a per cent of a star's
light at $300$\,K, habitable-zone hosts carry no mean warm excess at the level of
$10^{-5}$ of their light, no star within $15$\,pc brightened over a decade, and no star
that vanished from the Palomar sky reappears in the mid-infrared. The candidate
excesses that SPHEREx resolves are background sources, warm dust and cool companions.
The last of the companions' excesses, $3.3\sigma$ at $12$\um\ in HD~59099, is absent
from the WISE All-Sky release of the same exposures.

The limits apply to structures re-radiating as black bodies at the temperatures tested,
and to the epochs and sensitivities of the surveys: a $300$\,K enclosure of a star of
$0.01\,L_\odot$ within $210$\,pc moving faster than $19$ to $40$\,km\,s$^{-1}$ across
the line of sight, a $30$ to $100$\,K shell holding more than $0.2$ per cent of a solar
luminosity at $50$\,pc, and an enclosure at $200$ to $600$\,K of a vanished source
brighter than $R\approx19$. A sphere at a few tens of kelvin
around a star beyond $50$\,pc, one that emits mostly outside the WISE and AKARI bands,
a slow-moving enclosed star, or one whose waste heat is directed away from the Sun is
not constrained, and the SPHEREx fits do not reach components below $400$\,K. SPHEREx's
spectra of the dark companions and of the habitable-zone hosts would extend the
stacked test to $5$\um\ with resolved spectral shapes.

\section*{Acknowledgements}

This work has made use of data from the European Space Agency mission Gaia, processed by
the Gaia Data Processing and Analysis Consortium (DPAC), funded by national institutions,
in particular those participating in the Gaia Multilateral Agreement; of the Wide-field
Infrared Survey Explorer and NEOWISE, projects of the Jet Propulsion Laboratory,
California Institute of Technology, funded by NASA; of SPHEREx data from the NASA/IPAC
Infrared Science Archive; of AKARI, a JAXA project with the participation of ESA; of
the NASA Exoplanet Archive; of the Legacy Surveys; of the Spanish Virtual Observatory's
archive of vanishing objects; and of the SIMBAD and VizieR services at CDS, Strasbourg.

\begin{thebibliography}{99}\small\raggedright
\setlength{\itemsep}{1pt}

\bibitem{Dyson1960}
F.~J. Dyson, Search for artificial stellar sources of infrared radiation, \emph{Science}
\textbf{131}, 1667 (1960).

\bibitem{Carrigan2009}
R.~A. Carrigan, IRAS-based whole-sky upper limit on Dyson spheres, \emph{Astrophys. J.}
\textbf{698}, 2075 (2009).

\bibitem{Wright2014}
J.~T. Wright, B.~Mullan, S.~Sigurdsson and M.~S. Povich, The $\hat{G}$ infrared search
for extraterrestrial civilizations with large energy supplies. I. Background and
justification, \emph{Astrophys. J.} \textbf{792}, 26 (2014).

\bibitem{Griffith2015}
R.~L. Griffith, J.~T. Wright, J.~Maldonado, M.~S. Povich, S.~Sigurdsson and B.~Mullan,
The $\hat{G}$ infrared search for extraterrestrial civilizations with large energy
supplies. III. The reddest extended sources in WISE, \emph{Astrophys. J. Suppl.}
\textbf{217}, 25 (2015).

\bibitem{Suazo2022}
M.~Suazo, E.~Zackrisson, J.~T. Wright, et al., Project Hephaistos -- I. Upper limits on
partial Dyson spheres in the Milky Way, \emph{Mon. Not. R. Astron. Soc.} \textbf{512},
2988 (2022).

\bibitem{Suazo2024}
M.~Suazo, E.~Zackrisson, P.~K. Mahto, et al., Project Hephaistos -- II. Dyson sphere
candidates from Gaia DR3, 2MASS, and WISE, \emph{Mon. Not. R. Astron. Soc.}
\textbf{531}, 695 (2024).

\bibitem{Ren2024}
T.~Ren, M.~A. Garrett and A.~P.~V. Siemion, Background contamination of the Project
Hephaistos Dyson sphere candidates, \emph{Res. Notes AAS} \textbf{8}, 145 (2024).

\bibitem{Ren2026}
T.~Ren, M.~A. Garrett, E.~Zackrisson, et al., Archival diagnostics for potential
background contaminants of Project Hephaistos Dyson sphere candidates, \emph{Mon. Not.
R. Astron. Soc.}, in press, arXiv:2607.03619 (2026).

\bibitem{Korn2026}
A.~J. Korn, M.~Suazo, E.~Zackrisson, et al., Project Hephaistos -- III. Characterizing
anomalous infrared sources identified as Dyson-sphere candidates, arXiv:2607.25701
(2026).

\bibitem{Zackrisson2026}
E.~Zackrisson, A.~Bik, A.~Ali Asgar, et al., Project Hephaistos -- IV. James Webb Space
Telescope observations of two Dyson sphere candidates, arXiv:2607.09460 (2026).

\bibitem{Contardo2024}
G.~Contardo and D.~W. Hogg, A data-driven search for mid-infrared excesses among five
million main-sequence FGK stars, \emph{Astron. J.} \textbf{168}, 157 (2024).

\bibitem{Arenou2023}
Gaia Collaboration, F.~Arenou, C.~Babusiaux, et al., Gaia Data Release 3: stellar
multiplicity, a teaser for the hidden treasure, \emph{Astron. Astrophys.} \textbf{674},
A34 (2023).

\bibitem{Wright2010}
E.~L. Wright, P.~R.~M. Eisenhardt, A.~K. Mainzer, et al., The Wide-field Infrared Survey
Explorer (WISE): mission description and initial on-orbit performance, \emph{Astron. J.}
\textbf{140}, 1868 (2010).

\bibitem{Kervella2022}
P.~Kervella, F.~Arenou and F.~Th\'evenin, Stellar and substellar companions from Gaia
EDR3: proper-motion anomaly and resolved common proper-motion pairs, \emph{Astron.
Astrophys.} \textbf{657}, A7 (2022).

\bibitem{Kopparapu2013}
R.~K. Kopparapu, R.~Ramirez, J.~F. Kasting, et al., Habitable zones around
main-sequence stars: new estimates, \emph{Astrophys. J.} \textbf{765}, 131 (2013).

\bibitem{Akeson2013}
R.~L. Akeson, X.~Chen, D.~Ciardi, et al., The NASA Exoplanet Archive: data and tools for
exoplanet research, \emph{Publ. Astron. Soc. Pac.} \textbf{125}, 989 (2013).

\bibitem{Marocco2021}
F.~Marocco, P.~R.~M. Eisenhardt, J.~W. Fowler, et al., The CatWISE2020 catalog,
\emph{Astrophys. J. Suppl.} \textbf{253}, 8 (2021).

\bibitem{Murakami2007}
H.~Murakami, H.~Baba, P.~Barthel, et al., The infrared astronomical mission AKARI,
\emph{Publ. Astron. Soc. Jpn.} \textbf{59}, S369 (2007).

\bibitem{Mainzer2014}
A.~Mainzer, J.~Bauer, R.~M. Cutri, et al., Initial performance of the NEOWISE
reactivation mission, \emph{Astrophys. J.} \textbf{792}, 30 (2014).

\bibitem{Solano2022}
E.~Solano, B.~Villarroel and C.~Rodrigo, Discovering vanishing objects in POSS I red
images using the Virtual Observatory, \emph{Mon. Not. R. Astron. Soc.} \textbf{515},
1380 (2022).

\bibitem{Villarroel2020}
B.~Villarroel, I.~Imaz, J.~Bergstedt, et al., Exploring the Gaia Data Release 2 for
vanishing sources, \emph{Astron. J.} \textbf{159}, 8 (2020).

\bibitem{Dey2019}
A.~Dey, D.~J. Schlegel, D.~Lang, et al., Overview of the DESI Legacy Imaging Surveys,
\emph{Astron. J.} \textbf{157}, 168 (2019).

\bibitem{Isaacson2017}
H.~Isaacson, A.~P.~V. Siemion, G.~W. Marcy, et al., The Breakthrough Listen search for
intelligent life: target selection of nearby stars and galaxies, \emph{Publ. Astron.
Soc. Pac.} \textbf{129}, 054501 (2017).

\bibitem{Dore2014}
O.~Dor\'e, J.~Bock, M.~Ashby, et al., Cosmology with the SPHEREX all-sky spectral
survey, arXiv:1412.4872 (2014).

\end{thebibliography}

\end{document}
